\honest{3 Levels of Rationality Verification}{Eliezer Yudkowsky}{2009}

I strongly suspect that there could be an art of rationality beyond the standard skills, and beyond what any one person knows. What a group can do about it, I say, depends ``overwhelmingly'' on how we verify ``our many amazing good ideas.'' \nb{The examples that follow show that verification matters. None of them weighs it against other obstacles, so ``overwhelmingly'' is asserted.}

I sort verification into three levels. The first is reputational. A martial-arts master who fights realistic duels against other schools is at least not a complete poseur. Schools that never compete end up with masters who think they have chi powers, and so, I say, do the schools of psychoanalysis.

The second is experimental: teach randomly assigned students by method A or method B, measure them all the same way, and run statistics. Happiness research began, ``more or less,'' when people found that the question ``On a scale of 1 to 10, how good do you feel right now?'' agreed with other measures, and that its answers are useful when averaged over a hundred people. \nb{Validation studies bear this out. Sandvik, Diener and Seidlitz (1993) concluded that ``conventional self-report instruments validly measure'' subjective well-being.}

The third is organizational: a test of individuals that is hard to game. Can the SAT be beaten by studying just for the SAT? \nb{A 2009 review found that coaching raises SAT scores only a little: 10 to 20 points on the math section and 5 to 10 on critical reading, each scored from 200 to 800.} Should colleges grant their own degrees? I consider it ``drop-dead obvious'' that they should not, and decline to go into it. Does a hedge fund with 20\% returns beat the market, or is it selling puts? \nb{Andrew Lo built a hypothetical fund that sold puts every month; on 1990s prices its record beat the market by a wide margin.} A test that can be gamed, I say, makes the whole field game it and lose its purpose. \nb{This is Campbell's law, stated by Donald Campbell in 1976.}

In an added paragraph I note that noisy measures still serve experiments with enough subjects, but judging one person needs a precise one.

Then I put the question to you: how do you verify rationality skills, at any of the three levels? I propose no measure myself. ``Brainstorm, I beg you,'' and if you cannot think of a good idea, ``come up with a stupid one.'' Think, and do it before you read the other comments. \nb{Some measures of this kind already existed, such as a 2007 index of adult decision-making competence whose higher scorers reported fewer bad outcomes in their lives. Later came Stanovich, West and Toplak's Comprehensive Assessment of Rational Thinking (2016), and a 2015 study by the Center for Applied Rationality that found no change on its measures of cognitive biases after its workshops, in a study its author judged underpowered for those measures.}
